\begin{myChapterIntro}{本章内容}
\begin{itemize}
\item 对大型语言模型而言，“推理”意味着什么
\item 回顾 LLM 传统的预训练与后训练阶段
\item 介绍提升 LLM 推理能力的关键方法
\item 推理与模式匹配有何不同
\item 为什么要从零构建推理模型
\end{itemize}
\end{myChapterIntro}

欢迎来到大型语言模型（LLM）的下一个阶段：\emph{推理}。LLM 已经改变了我们处理和生成文本的方式，但它们的成功在很大程度上依赖统计模式识别。推理方法的新进展如今使 LLM 能够处理更复杂的任务，例如求解逻辑谜题和多步数学问题。

此外，推理是让 AI 智能体变得实用的关键技术，例如当智能体需要把任务拆解成步骤、使用工具并从错误中恢复时。推理 LLM 已被用于 OpenClaw 等智能体应用中。

本书将通过动手实践、代码优先的方式，揭示 LLM 推理方法的内部工作机制。我们会从一个预训练的 LLM 出发，逐步为它扩展推理能力。我们将从零实现这些推理组件，以了解这些方法在实践中如何运作。

如果想了解 LLM 本身是如何构建和训练的，我此前出版的、同样由 Manning 出版的\emph{《从零构建大语言模型》}（\href{http://mng.bz/orYv}{http://mng.bz/orYv}）详细介绍了这些基础知识，不过阅读本书并不需要先读那本书。

即使只对 LLM 的构建方式有粗略了解，读完本书后，也能理解推理模型的工作原理，并具备设计、制作原型和评估提升 LLM 推理能力主要方法的能力。

使用推理模型时，关键在于掌握如何自动检查数学等任务中的答案是否正确，以及如何把通用模型转变成更小的推理模型。本书通篇都会使用数学示例，因为它们易于自动检查，但跟进本书并不需要高深的数学知识。理解推理方法才是本书的核心重点。

本书注重实际应用与讲解，适合 LLM 工程师、机器学习研究者、应用科学家和软件开发人员阅读。理解推理模型是如何构建的，有助于我们判断什么时候标准 LLM 已经够用、什么时候额外的推理机制值得付出增加的复杂度、成本和延迟。这还有助于我们设计更好的评估方案、调试多步任务中的失败案例，并把推理技术应用到实际产品和研究工作流中。

\section{在 LLM 语境下定义推理}

本书中，我是在实际工程意义而非哲学意义上使用\emph{推理}一词的。就 LLM 而言，推理指的是模型在给出最终答案之前先生成中间步骤。

这些中间步骤可能表现为用户可见的逐步解释，也可能被包裹在诸如 \texttt{\textless{}think\textgreater{}...\textless{}/think\textgreater{}} 之类的特殊标签中、对用户隐藏。无论哪种情况，核心思想都一样：模型被引导或训练成用词元来承载中间解题步骤，而不是直接跳到结论。

相应地，我用\emph{推理模型}指代这样一类 LLM：它们通过训练或提示技术得到改进，能够生成这类中间步骤，从而提高在编程、逻辑谜题和数学问题等复杂任务上的准确率。

在进入本书的编码部分之前，我会先简要讨论改善这些推理行为的主要技术，以及它们与模式匹配和逻辑推理的关系。这将为后续讨论奠定基础：LLM 目前是如何构建的，它们如何处理推理任务，以及它们擅长什么、不擅长什么。

\begin{mySidebar}{思维链（chain-of-thought, CoT）}

这种生成中间步骤的风格通常被称为\emph{思维链}（CoT）。研究者和工程师常说模型在逐步“思考”问题，意思是它把中间推理过程显式地表达出来，让人更容易跟上。本书中，我是在上述常见的工程意义上使用\emph{推理}和\emph{思考}这两个词的。这并不意味着 LLM 真的会以和人类相同的方式推理或思考。

\end{mySidebar}

图 1.1 展示了传统 LLM 如何生成用户问题的答案。传统 LLM 可能不会展示它是如何得出答案的，而即使答案正确，也无助于用户理解背后的过程。

\myGraphic{0.92}{images/CH01_F01_Raschka2.png}{图 1.1 传统非推理 LLM 如何用一个简短答案回答问题的简化示意}

图 1.2 展示了 LLM 中进行多步（CoT）推理的一个简单示例。LLM 生成的中间推理步骤看起来很像一个人把内心的想法大声说出来。这些方法（以及由此产生的推理过程）在多大程度上贴近人类推理，仍是一个开放的研究问题，本书并不试图回答它。

\myGraphic{0.92}{images/CH01_F02_Raschka2.png}{图 1.2 推理 LLM 如何利用思维链处理多步推理任务的简化示意。模型并非只是回忆某个事实，而是把若干中间推理步骤组合起来，得出正确结论。中间推理步骤是否展示给用户，取决于具体实现。}

尽管图 1.2 是思维链推理的典型示例，但有必要强调：LLM 推理不同于传统的确定性推理。例如，符号逻辑引擎或定理证明器遵循严格的、基于规则的步骤，能保证一致性和正确性。它就像照着菜谱操作：每一步都是固定的，每次都必须产生相同的结果。只要严格按步骤操作，最终得到的一定是同一道菜。

与符号逻辑引擎不同，LLM 以\emph{自回归}方式生成推理，也就是根据训练数据中的统计模式一次预测一个词元。因此，LLM 的“推理步骤”并不能保证在逻辑上严谨，即使它们看上去很有说服力。

本书着重讲解并实现那些改善基于 LLM 的推理、让 LLM 更擅长处理复杂任务的基础技术。我希望，通过动手实践这些方法，读者能更好地理解和改进正在发展中的推理方法，甚至或许去探索它们与人类推理的异同。

\begin{mySidebar}{LLM 与人类推理}

LLM 中的推理过程可能表面上类似人类思维，尤其是在中间步骤的表述方式上。但重要的是要认识到一个关键区别：人类可以通过刻意运用逻辑规则，或基于对世界的内在模型进行推理，从而展开确定性推理。这里的“确定性”是指：如果从相同的事实出发、遵循相同的步骤，我们总会得出相同的结论。相比之下，当前的 LLM 推理是概率性的，也就是说，它根据训练数据中的统计模式一次生成一个词元，并不保证逻辑上的一致性。

人类常常通过有意识地操作概念、直觉地理解抽象关系，或从少量示例中归纳来推理。相比之下，LLM 是从海量文本中捕捉模式，而非依赖内置的推理规则或任何形式的有意识思考。

简言之，尽管经过推理增强的 LLM 的输出可能看起来很像人类，但其底层机制差别很大，并且仍是一个活跃的探索领域。

\end{mySidebar}

\section{理解标准的 LLM 训练流水线}

我们先简要看看传统（非推理）LLM 通常是如何训练的，从而理解其局限性所在。这些背景知识也有助于为接下来关于模式匹配与逻辑推理差异的讨论搭建框架。

在采用任何具体的推理方法之前，传统的 LLM 训练通常分为两个阶段：\emph{预训练}和\emph{后训练}，如图 1.3 所示。近期一些论文还在这两者之间划分出一个\emph{中期训练}（mid-training）阶段，例如在一个模型上继续用代码、数学或长上下文数据训练，但为了在本书中简化术语，我会把该阶段归入预训练。

\myGraphic{0.92}{images/CH01_F03_Raschka2.png}{图 1.3 典型 LLM 训练流水线概览。流程从一个权重随机初始化的模型开始，随后在大规模文本数据上进行预训练，通过预测下一个词元来学习语言模式。接着，后训练通过指令微调和偏好微调来改进模型，使 LLM 能更好地遵循人类指令并对齐人类偏好。}

在典型 LLM 训练流水线的预训练阶段，LLM 会在海量无标注文本上训练，数据量往往达数 TB，即数万亿个词元，来源包括书籍、网站、研究文章以及许多其他渠道。预训练的目标是让 LLM 学会预测这些文本中的下一个词（也就是下一个\emph{词元}）。

\begin{mySidebar}{词与词元}

词元是语言模型处理的文本最小单元。词元可以是一个完整的词、词的一部分，甚至标点符号，具体取决于分词器如何切分文本。

例如，“An LLM can be useful.”这句话可能会被常见的分词器切分成“An”、“ L”、“LM”、“ can”、“ be”、“ useful”和“.”这样的词元。随后这些词元会被转换成模型可以读取的数值 ID。

分词器并不属于 LLM 本身，但它仍是 LLM 文本处理与生成流水线中的关键组件。下一章我们会看到分词是如何工作的。

\end{mySidebar}

LLM 在海量数据集上预训练后会变得非常强大，这些数据集通常包含数 TB 的文本（相当于数万亿个词元）。这样的训练需要数千块 GPU 运行数月，成本可能高达数百万美元。这里的“强大”意味着 LLM 开始生成与人类写作高度相似的文本。在某种程度上，预训练的 LLM 还会开始展现出所谓的\emph{涌现能力}（emergent properties），即它们能够完成并未被明确训练过的任务，包括翻译和代码生成。

这些预训练模型作为基座模型，供后训练阶段使用。后训练采用两项关键技术：\emph{监督微调}（文献中常缩写为 \emph{SFT}，也叫\emph{指令微调}）和\emph{偏好调优}（preference tuning），用来教会 LLM 回应用户的查询。

如图 1.4 所示，指令微调提升了 LLM 在个人助理类任务上的能力，例如问答、摘要、文本翻译等。随后，偏好调优阶段会进一步改进这些能力。顾名思义，偏好调优有助于让回复贴合用户的偏好。有些读者可能熟悉基于人类反馈的强化学习（RLHF）这一术语，它是实现偏好调优的一种具体技术。

\myGraphic{0.92}{images/CH01_F04_Raschka2.png}{图 1.4 语言模型在不同训练阶段的回复示例。提示要求对睡眠与健康之间的关系做一段摘要。预训练的 LLM 给出的回答相关但不够聚焦，也没有直接遵循指令。经指令微调的 LLM 生成了与提示相符的简洁、准确的摘要。经偏好调优的 LLM 通过友善的语气和富有感染力的措辞进一步改进了回复，使答案更易引起共鸣、更以用户为中心。}

简言之，可以把预训练理解为“原始语言预测”（通过预测下一个词元），它赋予 LLM 一些能够生成连贯文本的基本属性和能力。后训练阶段则通过指令微调提升 LLM 对任务的理解，并通过偏好调优改进 LLM，使其给出符合偏好的风格化答案。

值得注意的是，即使是经过指令微调的模型也还不是“聊天机器人”。聊天界面增加了一层机制，在交互式多轮对话中引导模型的回复。这通常涉及系统提示、对话历史管理以及其他编排工作（附录 G 中实现了一个示例）。

\begin{myWarning}{注意}
 这些预训练和后训练阶段在我由 Manning 出版的\emph{《从零构建大语言模型》}（\href{http://mng.bz/orYv}{http://mng.bz/orYv}）一书中有所介绍。现在阅读的这本书并不要求详细了解这些阶段。这里，我们将加载一个已经历过昂贵的预训练和后训练阶段的模型，以便专注于推理模型特有的方法论。\end{myWarning}

\section{通过训练和推理技术提升 LLM 推理能力}

就 LLM 而言，推理进入公众视野，始于 2024 年 9 月 12 日 OpenAI 在 ChatGPT 中发布 o1。在这篇发布文章中，OpenAI 表示：“我们开发了一系列新的人工智能模型，旨在在回复之前花更多时间思考”（\href{https://openai.com/index/introducing-openai-o1-preview/}{https://openai.com/index/introducing-openai-o1-preview/}）。此外，OpenAI 还写道：“当处理科学、编程、数学及类似领域的复杂问题时，这些增强的推理能力可能特别有用。”

几个月后，也就是 2025 年 1 月，DeepSeek 发布了 DeepSeek-R1 模型及技术报告，其中详细介绍了开发推理模型的训练方法（\href{https://arxiv.org/abs/2501.12948}{https://arxiv.org/abs/2501.12948}）。这件事引起了巨大轰动，因为该公司不仅免费公开了一个能与专有的 o1 模型竞争并超越其性能的模型，还分享了训练此类模型的蓝图。本书中，我们将通过从零实现类似方法，来看看开发推理模型所用的方法论是如何运作的。

开发并提升 LLM 推理能力的不同方法可以归为三大类，如图 1.5 所示。这些方法通常应用于已经历过传统预训练和后训练阶段（包括指令微调和偏好调优）的 LLM。

\myGraphic{0.92}{images/CH01_F05_Raschka2.png}{图 1.5 提升 LLM 推理能力常用的三种方法：推理计算扩展、强化学习和蒸馏。这些方法通常在传统训练阶段（初始模型训练、预训练，以及包含指令微调和偏好调优的后训练）之后应用，但推理技术也可以应用于预训练基座模型。}

在本书中，我们将更详细地考察这些常见方法：

\begin{itemize}
\item \emph{推理时计算扩展}（也常称为\emph{推理计算扩展}、\emph{测试时扩展}等）包括在推理时（用户向模型输入提示时）提升模型推理能力、而无需训练或修改底层模型权重的方法。其核心思想是用增加的计算资源换取更好的性能，即使模型参数固定不变，也能通过思维链推理和各种采样流程等技术变得更强。这将是第 4 章和第 5 章的重点。
\item \emph{强化学习}（RL）指的是通过鼓励模型采取能带来高奖励信号的动作来提升其推理能力的训练方法。这些奖励可以是宽泛的，比如任务成功或启发式评分，也可以是狭义定义且可验证的，比如数学问题或编程任务中的正确答案。\newline 推理时计算扩展无需修改模型即可提升推理性能，而与之不同，RL 会在训练过程中更新模型的权重。这使模型能根据从环境中获得的反馈，通过试错来学习和改进推理策略。我们将在第 6 章和第 7 章中更详细地探讨 RL。
\item \emph{蒸馏}是指把规模更大、能力更强的模型所学到的复杂推理模式迁移到更小或更高效的模型中。就 LLM 而言，这通常意味着使用由更大、更强的模型生成的高质量标注指令数据集进行监督微调（SFT）。在 LLM 文献中，这项技术通常称为\emph{知识蒸馏}，或简称为\emph{蒸馏}。需要指出的是，这与深度学习中的传统知识蒸馏略有不同：在传统做法中，较小的（“学生”）模型通常同时从较大的（“教师”）模型产生的输出和 logits 中学习。这一话题将在第 8 章进一步讨论。
\end{itemize}

\begin{mySidebar}{面向推理与偏好调优的强化学习}

在开发推理模型时，重要的是把强化学习（RL）方法同基于人类反馈的强化学习（RLHF）区分开来，后者在开发传统 LLM 时的偏好调优阶段使用。两种情形都使用相同的底层过程（RL），但主要区别在于奖励是如何获得和验证的（RLHF 依赖人类判断，而推理 RL 依赖自动验证器或环境）。

RLHF 将显式的人类评估或对模型输出的排序作为奖励信号，直接把模型引导向人类偏好的行为。相比之下，推理模型语境下的 RL 通常依赖自动化的或基于环境的奖励信号，这类信号可能更客观，但与人类偏好的一致性更低。例如，在推理模型开发流水线中，RL 可以通过对正确性给出显式奖励，来训练模型精通数学证明。相反，RLHF 则会让人类评估者对各种回复进行排序，以鼓励输出与人类标准和主观偏好高度一致。

\end{mySidebar}

\section{模式匹配与逻辑推理}

在预训练期间，LLM 会接触海量文本，并学会通过识别和重现数据中的统计关联来预测下一个词元。这一过程使它们能生成流畅、连贯的文本，但其基础本质上是学到的统计规律，而非显式规则或有保证的逻辑推理。

LLM 对提示作出回应时，会生成在统计上与训练期间所见模式一致的文本续写。这其中包含许多从前提推及结论的模式，看起来像是逻辑推理。关键区别在于，模型并不显式地表示前提并应用形式化规则，而是根据训练数据中的相似示例预测可能的续写。

来看下面这个例子：

\begin{mySidebarBox}
\textbf{\includegraphics[width=0.05\linewidth, valign=c]{images/Prompt-Icon.png}} 德国的首都是……
\end{mySidebarBox}

\begin{mySidebarBox}
\textbf{\includegraphics[width=0.05\linewidth, valign=c]{images/Response-Chatgpt.png}} 柏林
\end{mySidebarBox}

LLM 给出“柏林”这个答案，并不是在逻辑地推导。它实际上是在回忆从训练数据中学到的一种强统计关联。这种行为正是我们所说的\emph{模式匹配}，也就是说，模型是依据学到的相关性来完成文本的，而不是通过应用结构化的推理步骤。

但对于超出模式识别范围的任务、那些正确答案取决于从给定事实中得出结论的任务，又当如何呢？这就引出了另一种不同的能力：逻辑推理。

逻辑推理指运用规则系统性地得出结论。与模式匹配不同，它依赖中间推理步骤，以及识别矛盾、依据形式化关系进行推演的能力。

以下面这个提示为例：

\begin{mySidebarBox}

\textbf{\includegraphics[width=0.05\linewidth, valign=c]{images/Prompt-Icon.png}} 所有鸟都会飞。企鹅是鸟。企鹅会飞吗？

\end{mySidebarBox}

对此有两种评判方式。

第一种，在\emph{封闭世界}（仅依据提示）的设定下，给定提示中的两个前提（主张或假设），有效的答案是“是的，企鹅会飞”。

第二种，在\emph{开放世界}（带背景知识）的设定下，如果我们允许引入提示中未包含的知识（比如企鹅不会飞），那么这个外部事实就与从前提推出的结论相冲突，如图 1.6 所示。推理系统应当注意到这种不一致，并要么请求澄清，要么弱化第一个陈述（例如“大多数鸟会飞，但有企鹅这样的例外。”）。

\myGraphic{0.92}{images/CH01_F06_Raschka2.png}{图 1.6 相互矛盾的前提导致逻辑不一致。由“所有鸟都会飞”和“企鹅是鸟”，我们推出“企鹅会飞”。这一结论与既有事实“企鹅不会飞”相冲突，从而产生矛盾。}

相比之下，统计式（模式匹配）的 LLM \emph{不会}显式地追踪矛盾，而是依据学到的文本分布进行预测。例如，如果“所有鸟都会飞”这类信息在训练数据中被强烈强化，模型可能会自信地回答：“是的，企鹅会飞。”

下一节，我们会看一个 LLM 如何处理这个“所有鸟都会飞……”提示的例子。

\begin{mySidebar}{逻辑推理与基于规则的系统}

为什么显式的基于规则的系统不那么流行？基于规则的系统在 20 世纪 80 年代和 90 年代被广泛用于医学诊断、法律判决和工程领域。它们在医疗、法律、航空航天等关键领域仍在使用，这些领域往往要求可解释的推理和透明的决策过程。这类系统难以实现，因为它们很大程度上依赖人工编写的启发式规则。（启发式规则是一些简单的决策规则，能快速给出一个足够好的答案，但不保证最优。）相比之下，包括 LLM 在内的深度神经网络并不执行手写规则；它们从数据中学习决策模式，在大规模训练时可以非常灵活。

\end{mySidebar}

\section{不依赖显式规则来模拟推理}

我们已经看到矛盾的前提如何导致逻辑不一致。传统 LLM 不会显式追踪矛盾，而是依据学到的文本分布生成回复。

图 1.7 展示了像 OpenAI ChatGPT 中的 GPT-4o 这类未经推理增强的 LLM，如何回应上一节讨论的“所有鸟都会飞……”提示。可以看到，尽管 GPT-4o 并不被视为推理模型（不同于 OpenAI 的其他产品，比如 o1、o3、o4-mini，以及更近期的 GPT-5，这些都是明确采用推理方法论开发的），它似乎也能正确作答。

\myGraphic{0.92}{images/CH01_F07_Raschka2.png}{图 1.7 语言模型（ChatGPT 中的 GPT-4o）如何看似对矛盾前提进行“推理”的示例}

那么，4o 模型是如何生成答案的？这是否意味着 GPT-4o 在显式地进行逻辑推理？不，未必如此。至少可以说，GPT-4o 在熟悉的情境中非常擅长模拟逻辑推理。

GPT-4o 并未实现显式的矛盾检查，而是依据按概率加权的模式生成答案。如果训练数据中包含大量纠正该矛盾的实例（例如“企鹅不会飞”这样的文本），使模型学到“企鹅”与“不会飞”之间的统计关联，那么这种做法就足够奏效。换言之，模型之所以能隐式地识别出矛盾，是因为它在训练中频繁遇到过这一确切的推理场景。这种效果在很大程度上依赖训练数据中大量类似推理模式所建立的统计关联。因此，即使传统 LLM 看起来在进行逻辑演绎，如图 1.7 所示，它也不是在执行显式的、基于规则的逻辑，而是在利用其海量训练数据中的模式。

尽管如此，GPT-4o 在此处的成功很好地说明，隐式模式匹配在大规模训练后能变得多么强大。这类基于模式的推理模型通常在两种情形下会表现不佳：

\begin{itemize}
\item 当逻辑场景是全新的（训练数据中未曾出现过）时
\item 当推理复杂度很高、涉及错综复杂的多步逻辑关系时
\end{itemize}

\begin{mySidebar}{逻辑推理与当前已有的推理模型}

尽管 GPT-4o 并未被正式归类为推理模型，但 OpenAI 提供了若干专用的推理模型，包括 o1、o3、o4-mini 和 GPT-5。此外，其他公司也在开发具备显式推理能力的 LLM。截至本书写作时，常被提及的例子包括 Anthropic 的 Claude 4、xAI 的 Grok 4、Google 的 Gemini 2.5、DeepSeek 的 R1、阿里巴巴的 Qwen3 等等。这些模型所采用的技术正是本书的重点。正如我们将看到的，这一切的实现并不需要基于规则的推理流水线。相反，LLM 会通过改进后的推理与训练方法论来学习或提升其推理能力。

\end{mySidebar}

或许可以说，LLM 通过学到的模式来模拟逻辑推理。我们可以借助特定的推理方法进一步改进这一点，包括推理计算扩展和强化学习等后训练策略，但这些 LLM 内部并没有显式执行任何基于规则的逻辑。

此外，LLM 处于一个连续谱系上。早在 OpenAI 的 o1 和 DeepSeek-R1 等专用推理模型出现之前，LLM 就已经能够模拟推理行为。这些模型表现出的行为符合本章 1.1 节对推理的早期定义，即在给出最终答案之前先生成中间步骤。我们现在明确贴上的“推理模型”标签，本质上是这一能力更精细的版本。这些得到提升的推理能力，是通过特定的推理计算扩展技术（将在第 4 章和第 5 章讨论）以及有针对性的后训练方法（如第 6 章和第 7 章的强化学习）实现的，后者的目的是改进并强化类推理行为。

\section{为什么要从零构建推理模型？}

自 2025 年 1 月 DeepSeek-R1 发布以来，提升 LLM 的推理能力已成为 AI 领域最热门的话题之一，这也不无道理。更强的推理能力让 LLM 能够处理更复杂的问题，使其在用户关心的各类任务上表现更出色。

这一转变也体现在 OpenAI 首席执行官 2025 年 2 月 12 日的一段声明中：

上面这段话凸显了领先的 LLM 提供商向推理模型转变的重大趋势。这里，“思维链”指的是一种引导语言模型逐步推理以提升其推理能力的提示技术；我们将在第 4 章和第 5 章更详细地介绍它。

同样值得注意的是其中提到要知道“何时该长时间思考，何时不必”。这暗示了一个重要的设计考量：推理并非总是必要或可取的。例如，推理模型是为求解谜题、高等数学问题以及高难度编程任务等复杂任务而设计的。它们在智能体系统中也常常很有用，尤其是在任务分解、规划以及选择使用哪些工具方面。对于摘要、翻译或基于知识的问答等较简单的任务，它们并非必需。事实上，凡事都用推理模型可能既低效又昂贵。例如，它们通常使用成本更高、输出更啰嗦，有时还会因“过度思考”而更容易出错。这里适用一条简单的法则：为任务选择恰当的工具（或恰当类型的 LLM）。

推理模型往往比非推理模型更昂贵，原因有两个。其一，它们倾向于产生更长的输出，因为其中包含解释答案如何得出的中间步骤。如图 1.8 所示，LLM 一次生成一个词元，而每个词元都需要一次\emph{完整的前向传播}。如果推理模型的答案长度翻倍，生成过程就大致需要两倍的前向传播次数，从而增加计算成本。

\myGraphic{0.92}{images/CH01_F08_Raschka2.png}{图 1.8 LLM 中逐个词元的生成过程。每一步中，LLM 会接收目前为止生成的完整序列，并预测下一个词元，该词元可能代表一个词、子词或标点符号，具体取决于分词器。新生成的词元会被追加到序列末尾，并作为下一步的输入。这种迭代式解码过程既用于标准语言模型，也用于以推理为核心的模型。}

其二，许多推理工作流需要针对单个任务多次运行模型，例如采样多个候选解法、调用工具或运行验证器。这些额外的调用会使处理的词元总数成倍增加，从而在图 1.8 所示的单次调用行为之外进一步推高成本。

这正是从头实现这些模型和方法大有裨益的原因。这是理解它们如何运作的最佳方式之一。而如果我们理解了 LLM 和推理模型如何运作，就能更好地理解这些权衡。

\section{从零构建推理模型的路线图}

既然我们已经从宏观角度讨论了 LLM 中的推理，接下来的各章将带你从零开始编写并应用推理方法。从高层次看，路线图很简单。我们会从一个传统 LLM 出发，学习如何评估其推理行为，然后探索改进它的两大类方法：推理技术和训练技术。图 1.9 展示了主要步骤。

\myGraphic{0.92}{images/CH01_F09_Raschka2.png}{图 1.9 本书将构建内容的高层路线图。我们会从一个传统 LLM 出发，加入评估方法以便衡量其进展，然后探索改进推理的两大类方法：推理技术和训练技术。}

在图 1.9 中，阶段 1 是一个明确的步骤，因为在实践中，推理方法通常应用于已有的基座 LLM，而不是从头训练的模型。我们会从一个已经预训练好的传统 LLM 出发，把它当作我们希望改进的基线。

阶段 2 实现所需的评估工具，用来判断某个方法是否真的有帮助。正因如此，阶段 3 和阶段 4 都会回到评估环节：在通过推理时方法或额外训练改进推理行为之后，我们回到阶段 2，衡量这些改动是否奏效。

阶段 3 在不改变模型权重的情况下改进推理时的推理行为，而阶段 4 通过额外训练改进模型本身，把预训练的 LLM 变成一个专用的推理模型。

图 1.10 提供了更多细节，并概括了各章涵盖的不同子步骤。

\myGraphic{0.92}{images/CH01_F10_Raschka2.png}{图 1.10 章级子步骤的详细路线图。加载基座模型之后，我们会介绍基于基准测试和基于评判的评估，然后是高级文本生成、投票和自我修正等推理时方法，最后是基于强化学习和蒸馏的训练时方法。}

我很期待接下来的旅程，也希望各位同样如此。

\section{小结}

\begin{itemize}
\item 传统的 LLM 训练分为若干阶段： 预训练：模型从海量文本中学习语言模式 指令微调：改进模型对用户提示的回复 偏好调优：使模型输出与人类偏好对齐
\item 推理方法是在传统 LLM 之上应用的。
\item LLM 中的推理指的是改进模型，使其在给出最终答案之前显式生成中间步骤（思维链），这往往能提高多步任务的准确率。
\item LLM 中的推理不同于基于规则的推理，其运作方式很可能也不同于人类推理。目前，主流共识是 LLM 中的推理依赖统计模式匹配。
\item LLM 中的模式匹配完全依赖从数据中学到的统计关联，能够生成流畅的文本，但缺乏显式的逻辑推理。
\item 改进 LLM 推理能力可以通过几种方式实现： 推理时计算扩展无需重新训练即可改进推理（例如思维链提示）。 强化学习通过奖励信号对模型进行显式训练。 监督微调和蒸馏，使用来自更强推理模型的示例。
\item 从零构建推理模型，能让人切实了解 LLM 的能力、局限以及计算上的权衡。
\end{itemize}
